\documentclass[UTF8,zihao=-4]{ctexart}
\usepackage[a4paper,margin=2.5cm]{geometry}
\usepackage{amsmath,amssymb,booktabs,graphicx,float}
\usepackage{caption}
\usepackage[hidelinks]{hyperref}
\usepackage{enumitem}
\graphicspath{{../../figures/}{../../../results/figures/}}
\setlength{\parindent}{2em}
\setlength{\parskip}{0pt}
\pagestyle{plain}

\title{B 题论文标题（不得出现身份信息）}
\author{}
\date{}

\begin{document}

% 电子版第一页必须为摘要专用页；不要加入承诺书和编号专用页。
\maketitle
\vspace{-2.5em}
\begin{abstract}
在此撰写摘要。依次交代问题、方法、关键量化结果和结论，原则上不超过一页。

\textbf{关键词：} 关键词一；关键词二；关键词三
\end{abstract}
\thispagestyle{plain}
\clearpage

\section{问题重述}
用自己的语言准确重述问题，明确输入、输出、约束与目标，避免大段照抄题面。


\section{问题分析}
说明各子问题之间的关系、解题路径、模型选择依据和可验证指标。


\section{模型假设与符号说明}
逐条给出必要、合理且可检验的假设，并说明主要符号、单位和取值范围。


\section{模型建立与求解}
按“建立模型—算法实现—参数设置—求解结果”的顺序展开。每个结果都应能由支撑材料复现。


\section{模型检验、敏感性与误差分析}
进行模型检验、基线对比、敏感性分析、稳健性分析或误差分析，说明适用边界。


\section{模型评价与推广}
客观总结模型优点、局限、改进方向与可推广场景。


\section{结论}
逐一回应题目要求，给出明确、量化、可解释的最终结论。


\begin{thebibliography}{99}
\bibitem{example} 作者. 文献标题[J]. 期刊, 年, 卷(期): 页码.
\end{thebibliography}

\appendix
\section{支撑材料文件清单}
\begin{itemize}[leftmargin=2em]
  \item \texttt{src/}：完整可运行源程序；
  \item \texttt{data/}：必要的自主查阅与处理后数据；
  \item \texttt{results/}：必要的中间结果与表格。
\end{itemize}

\section{补充推导或结果}
在此放置必要的补充内容。

\end{document}
